\chapter{Fainting (Syncope)}

\section*{Definition}
Syncope (fainting) is a transient loss of consciousness due to global cerebral hypoperfusion, characterized by rapid onset, brief duration, and spontaneous complete recovery. It is a common symptom with many potential causes.

\section*{When to See a Doctor}
\begin{itemize}
\item Seek emergency care for fainting during exercise, chest pain, palpitations, severe shortness of breath, major injury, prolonged confusion, convulsions, new neurologic symptoms, or a family history of sudden cardiac death.
\item Arrange medical assessment after a first faint, recurrent fainting, fainting without a clear trigger, or fainting with new symptoms.
\end{itemize}

\section*{How It Develops}
Cerebral blood flow requires adequate cardiac output, systemic vascular resistance, and cerebral perfusion pressure. Syncope occurs when there is a sudden reduction in one or more of these factors: reflex-mediated bradycardia and widening of blood vessels (reflex syncope), orthostatic hypotension with inadequate compensation, or cardiac causes with reduced pump function or arrhythmias.

\section*{Causes}
\begin{itemize}
\item Vasovagal syncope (emotional stress, pain, fear, prolonged standing)
\item Orthostatic hypotension (dehydration, blood loss, medications)
\item Situational syncope (micturition, defecation, coughing, swallowing)
\item Cardiac arrhythmias (bradycardia, heart block, ventricular tachycardia)
\item Structural heart disease (aortic stenosis, hypertrophic cardiomyopathy)
\item Pulmonary embolism
\item Medication-induced (vasodilators, diuretics, beta-blockers)
\end{itemize}

\section*{Related Symptoms}
\begin{itemize}
\item Dizziness or lightheadedness
\item Presyncope (near-fainting)
\item Palpitations
\item Seizures (may be confused with syncope)
\item Fall or injury from loss of consciousness
\end{itemize}

\section*{Medical Evaluation}
\begin{itemize}
\item Detailed history of the event and triggers
\item Orthostatic vital signs (blood pressure and heart rate supine and standing)
\item 12-lead electrocardiogram to exclude arrhythmia
\item Echocardiography only when history, examination, or ECG suggests structural heart disease
\item Holter or event monitor for suspected arrhythmia
\item Tilt-table test for recurrent vasovagal syncope
\end{itemize}

\section*{Treatment Options}
\begin{itemize}
\item Reassurance and education for reflex syncope
\item For selected reflex or orthostatic syncope, individualized fluid and salt advice may help when not contraindicated
\item Compression stockings for orthostatic hypotension
\item Discontinue or adjust offending medications
\item Pacemaker for bradycardia or heart block
\item ICD for ventricular arrhythmias
\item Specific therapy for underlying cardiac or pulmonary cause
\end{itemize}

\section*{Self-Care and Home Management}
\begin{itemize}
\item If feeling faint, lie down or sit with head between knees
\item Cross legs and tense muscles to increase blood pressure
\item Stay well hydrated
\item Avoid prolonged standing in hot environments
\item Stand up slowly from sitting or lying position
\item Avoid alcohol and heavy meals before activities
\item Recognize and avoid individual triggers
\end{itemize}

\section*{References}
\begin{thebibliography}{99}
\bibitem{fainting:presyncope1} Shen WK, Sheldon RS, Benditt DG, et al. ``2017 ACC/AHA/HRS guideline for the evaluation and management of patients with syncope.'' Circulation. 2017;136(5):e60-e122.
\bibitem{fainting:presyncope2} Brignole M, Moya A, de Lange FJ, et al. ``2018 ESC guidelines for the diagnosis and management of syncope.'' Eur Heart J. 2018;39(21):1883-1948.
\bibitem{fainting:presyncope3} Kapoor WN. ``Syncope.'' N Engl J Med. 2000;343(25):1856-1862.
\bibitem{fainting:presyncope4} Albassam OT, Redelmeier RJ, Shadowitz S, et al. ``Did this patient have cardiac syncope?'' JAMA. 2019;321(24):2448-2457.
\bibitem{fainting:presyncope5} Serrano LA, Hess EP, Bellolio MF, et al. ``Accuracy and quality of clinical decision rules for syncope in the emergency department.'' Acad Emerg Med. 2010;17(8):799-809.
\bibitem{fainting:presyncope6} Kaufmann H, Bhatt DG. ``Syncope and presyncope.'' In: Harrison\'s Principles of Internal Medicine. 21st ed. McGraw-Hill; 2022.
\end{thebibliography}
